\documentclass[11pt,fontset=fandol]{ctexart}

\usepackage[a4paper,margin=2.5cm]{geometry}
\usepackage{amsmath,amssymb,mathtools}
\usepackage{xcolor}
\usepackage{hyperref}
\usepackage{microtype}
\usepackage{enumitem}
\usepackage{titlesec}

\definecolor{ink}{HTML}{1F2933}
\definecolor{muted}{HTML}{52606D}
\definecolor{accent}{HTML}{0F766E}
\definecolor{soft}{HTML}{F0Fdfa}

\hypersetup{
  colorlinks=true,
  linkcolor=accent,
  urlcolor=accent
}

\titleformat{\section}
  {\Large\bfseries\color{ink}}
  {\thesection}{0.8em}{}

\setlist[itemize]{leftmargin=1.5em}
\setlength{\parindent}{0pt}
\setlength{\parskip}{0.7em}

\title{\bfseries Lipschitz 常数小为什么允许更大的学习率}
\author{}
\date{}

\begin{document}
\maketitle
\vspace{-2em}

\begin{center}
\colorbox{soft}{
  \parbox{0.9\linewidth}{
    \textbf{核心结论：}
    若目标函数的梯度是 \(L\)-Lipschitz 连续的，则梯度下降中保证单步下降的学习率通常满足
    \[
      0<\eta<\frac{2}{L}.
    \]
    因此，\(L\) 越小，允许的学习率上界越大。
  }
}
\end{center}

\section{问题设定}

考虑目标函数 \(f(x)\)，梯度下降的更新规则为
\[
  x_{t+1}=x_t-\eta \nabla f(x_t),
\]
其中 \(\eta\) 是学习率。

假设 \(f\) 的梯度是 \(L\)-Lipschitz 连续的，即
\[
  \|\nabla f(x)-\nabla f(y)\|\le L\|x-y\|.
\]
这意味着梯度变化不会超过输入变化的 \(L\) 倍。\(L\) 越小，梯度场越平滑。

\section{下降引理}

由梯度 Lipschitz 连续可得 \(f\) 是 \(L\)-smooth 函数，并满足下降引理：
\[
  f(y)\le f(x)+\nabla f(x)^\top (y-x)+\frac{L}{2}\|y-x\|^2.
\]

这个不等式说明：在点 \(x\) 附近，函数值可以被一个一阶项加二次误差项控制。二次误差项中的 \(L\) 决定了函数曲面的最大弯曲程度。

\section{代入梯度下降更新}

令
\[
  y=x-\eta \nabla f(x).
\]
代入下降引理：
\[
\begin{aligned}
  f(x-\eta \nabla f(x))
  &\le
  f(x)+\nabla f(x)^\top(-\eta \nabla f(x))
  +\frac{L}{2}\|-\eta \nabla f(x)\|^2 \\
  &=
  f(x)-\eta\|\nabla f(x)\|^2
  +\frac{L\eta^2}{2}\|\nabla f(x)\|^2 \\
  &=
  f(x)-\eta\left(1-\frac{L\eta}{2}\right)\|\nabla f(x)\|^2.
\end{aligned}
\]

\section{学习率条件}

为了保证目标函数下降，需要
\[
  \eta\left(1-\frac{L\eta}{2}\right)>0.
\]
由于 \(\eta>0\)，因此只需
\[
  1-\frac{L\eta}{2}>0.
\]
于是得到
\[
  \eta<\frac{2}{L}.
\]

因此，允许的最大学习率与 Lipschitz 常数 \(L\) 成反比：
\[
  \eta_{\max}\propto \frac{1}{L}.
\]

\section{更保守的常用选择}

实际分析中常用更保守的学习率范围：
\[
  0<\eta\le \frac{1}{L}.
\]
此时
\[
  1-\frac{L\eta}{2}\ge \frac{1}{2},
\]
因此有更清晰的下降量下界：
\[
  f(x-\eta\nabla f(x))
  \le
  f(x)-\frac{\eta}{2}\|\nabla f(x)\|^2.
\]

\section{直观解释}

\begin{itemize}
  \item \(L\) 小：梯度变化平缓，函数曲面较光滑，可以迈更大的步。
  \item \(L\) 大：梯度变化剧烈，函数曲面更陡或更弯，步长过大容易震荡或发散。
  \item 因此，Lipschitz 常数越小，理论上允许的稳定学习率越大。
\end{itemize}

\end{document}
